\documentclass[tikz,border=4pt]{standalone}
\usepackage{pe-tikz}
\begin{document}
\begin{tikzpicture}[
  tbl/.style={pe/box, text width=30mm, align=left, font=\scriptsize, inner sep=3pt},
  hd/.style={font=\small\bfseries, text=ink}]
  % PannonShop
  \node[hd, text=cB] at (-4.2,2.6) {PannonShop (üzleti eset)};
  \node[tbl, draw=cB] (t1) at (-5.9,1.1) {\textbf{Számlatételek}\\Invoice, StockCode,\\Quantity, InvoiceDate,\\Price, CustomerID};
  \node[tbl, draw=cB] (t2) at (-2.5,1.1) {\textbf{Termékek}\\kategória, alkategória,\\listaár};
  \node[tbl, draw=cB] (t3) at (-5.9,-0.8) {\textbf{Vevők}\\régió, regisztráció};
  \node[tbl, draw=cB] (t4) at (-2.5,-0.8) {\textbf{Értékelések}\\csillag, szöveg};
  \draw[pe/thin, cB] (t1) -- (t2); \draw[pe/thin, cB] (t1) -- (t3); \draw[pe/thin, cB] (t4) -- (t2);
  % Üzem
  \node[hd, text=cD] at (4.2,2.6) {Pannon üzem (mérnöki eset)};
  \node[tbl, draw=cD] (u1) at (2.5,1.1) {\textbf{Folyamatmérések}\\10 perces: hőmérséklet,\\áramok, nyomás, energia};
  \node[tbl, draw=cD] (u2) at (5.9,1.1) {\textbf{Laborminták}\\4 óránként: koncentráció,\\szennyezés, megfelelőség};
  \node[tbl, draw=cD] (u3) at (2.5,-0.8) {\textbf{Riasztásnapló}\\esemény, időbélyeg};
  \node[tbl, draw=cD] (u4) at (5.9,-0.8) {\textbf{Karbantartás}\\tisztítások, termékváltás};
  \draw[pe/thin, cD] (u1) -- (u2); \draw[pe/thin, cD] (u1) -- (u3); \draw[pe/thin, cD] (u1) -- (u4);
  % közös kérdések
  \node[pe/stepfill, text width=125mm, align=left, font=\scriptsize] (q) at (0,-3.0) {
    \textbf{Végigvonuló kérdések:}\enspace
    adattárház és OLAP (2.) \textbullet\ tisztítás (3.) \textbullet\ feltárás (4.) \textbullet\ szegmentálás, üzemállapotok (5.) \textbullet\
    kosárelemzés, riasztássorozatok (6.) \textbullet\ lemorzsolódás, minőség (7., 9.) \textbullet\ lágy szenzor (8., 10.) \textbullet\
    csalás és folyamathiba (11.) \textbullet\ előrejelzés (12.) \textbullet\ vélemények (13.) \textbullet\ összegző esettanulmányok (14., 15.)};
  \draw[pe/flow] (-4.2,-1.55) -- (-4.2,-2.45);
  \draw[pe/flow] (4.2,-1.55) -- (4.2,-2.45);
\end{tikzpicture}
\end{document}
